\section{Introduction}
Open-vocabulary scene models support language, reference-image and spatial
queries by connecting semantic evidence to a three-dimensional scene. Their
outputs must satisfy both identity and extent: locating the right object does
not ensure that its complete support is recovered. A response that agrees
better with one observation may omit occluded parts, merge distractors or
replace a correct answer unnecessarily. Preserving teacher features alone
therefore does not preserve a system's effective query abilities.

We investigate a focused question: \emph{given an existing complete answer and
alternative answers from fixed providers, which observable evidence justifies
replacement?} This moves the research variable to the final decision. We retain
native teacher matching, geometric carriers and complete decoders, and predict
relative task utility rather than interpreting teacher similarity as answer
quality. The incumbent is an explicit candidate; acceptance returns an intact
alternative rather than a patchwork of elements from different answers.

The proposed mechanism has three connected steps. It constructs complete
candidates with their continuous response and actual selection; summarizes
query association and regional overlap in teacher-order-invariant coordinates;
and applies one shared utility comparison with calibrated retention. We also
derive an exact finite-candidate decomposition of performance into available
gain and decision regret. This decomposition makes failures testable: inadequate
candidates, incorrect ranking and excessive or beneficial retention have
different implications for the next experiment.

Our contribution claims are tied to tests. The complete-answer formulation
establishes a common intervention without changing representations. The
query-conditioned decision must outperform controls that remove or permute
query correspondence. Selective retention must improve the repair--harm tradeoff
relative to the same model without a positive margin. These tests require no
feature-dimension, compression or combinatorial system ablations. The present
paper reports the completed comparisons and controls, while identifying the
remaining experiments needed to establish the central empirical claim.
